% 来源：sources/2016-2025-期中-甲-历年真题汇编.pdf 第 8 页
\begin{problem}{10 分}{9.5cm}
设
\[
D_n=\begin{vmatrix}
4&-3&0&\cdots&0&0\\
-1&4&-3&0&\cdots&0\\
0&-1&4&-3&0&\cdots\\
\vdots&\vdots&\vdots&\vdots&\vdots&\vdots\\
0&\cdots&0&-1&4&3\\
0&0&\cdots&0&-1&4
\end{vmatrix},
\]
试证明：$D_n=\dfrac{3^{n+1}-1}{2}$。
\end{problem}

\begin{problem}{10 分}{8cm}
设 $B=\begin{pmatrix}1&-5&1\\0&-1&0\\2&-2&3\end{pmatrix}$，$C=\begin{pmatrix}1&0&0\\0&2&0\\0&0&-1\end{pmatrix}$，三阶方阵满足方程 $2BA+CB=O$，其中 $O$ 为三阶零方阵，试求 $|A+3E|$。
\end{problem}

\begin{problem}{15 分}{8.5cm}
设有实方阵 $A=\begin{pmatrix}1&2&3\\2&4&t\\3&t&9\end{pmatrix}$，其中 $t$ 为实常数，已知 $B$ 为三阶非零方阵且满足 $BA=O$，其中 $O$ 为三阶零方阵，试求 $t$ 的值以及 $B$ 的秩的可能取值。
\end{problem}

\begin{problem}{15 分}{9.5cm}
设 $a,b,c$ 为实常数，且满足 $b^{2}\ne ac$，$a+b+c=0$，试证明线性方程组
\[
\begin{cases}
ax_1+bx_2+c=0,\\
bx_1+cx_2+a=0,\\
cx_1+bx_2+b=0
\end{cases}
\]
有唯一解并求出这个唯一解。
\end{problem}

\begin{problem}{15 分}{8.5cm}
设方阵 $A$ 的伴随矩阵为
$A^{*}=\begin{pmatrix}1&0&0&0\\0&1&0&0\\1&0&1&0\\0&-3&0&8\end{pmatrix}$，
已知矩阵 $B$ 满足 $AB=B+3A$，试求矩阵 $B$。
\end{problem}

\begin{problem}{20 分}{12cm}
设 $A=\begin{pmatrix}1&\cdots&1\\\vdots&\ddots&\vdots\\1&\cdots&1\end{pmatrix}_{n\times n}$。求
\begin{enumerate}[label=(\arabic*),leftmargin=2.2em,itemsep=0.25em,topsep=0.3em]
  \item 一个二次实系数多项式 $f(x)=x^{2}+ax$ 使得 $f(A)$ 为二阶零方阵；
  \item $A^{100}$；
  \item $(A+E)^{3}$；
  \item $(A+E)^{-1}$。
\end{enumerate}
\end{problem}

\begin{problem}{8 分}{6.5cm}
设 $A,B$ 是两个 $n$ 阶实方阵，试证明：$r(A)=r(AB)$ 当且仅当存在 $n$ 阶实方阵 $C$ 使得 $A=ABC$。
\end{problem}

\begin{problem}{7 分}{6cm}
设 $A$ 为 $n$ 阶方阵，试证明：存在对称矩阵 $S$ 和可逆矩阵 $P$ 使得 $A=SP$。
\end{problem}
